\section*{Chapitre \ref{ch:b2:funcseq} --- Suites et séries de
fonctions}

\begin{solution}{exo:b2:funcseq:1}
$f_n(x) = \frac{x}{1 + nx}$~: limite simple $0$ sur $\intcc{0}{1}$.
\omterm{def:b2:funcseq:def}{Uniformément}~: $f_n$ croît sur $\intcc{0}{1}$ (dérivée
$\frac{1}{(1+nx)^2} > 0$), donc $\norm{f_n}_\infty = f_n(1) =
\frac{1}{1+n} \to 0$~: convergence \emph{uniforme} sur $\intcc{0}{1}$.

$g_n(x) = nx\,\eu^{-nx^2}$~: limite simple $0$ (l'exponentielle
l'emporte). Sup~: $g_n' = n\eu^{-nx^2}(1 - 2nx^2)$ s'annule en $x_n =
\frac{1}{\sqrt{2n}}$, où $g_n(x_n) = \sqrt{\frac n2}\,\eu^{-1/2}
\to \infty$~: pas uniforme sur $\intcc{0}{1}$ --- mais uniforme sur
$\intco{\delta}{1}$, puisque là $g_n(x) \leq n\,\eu^{-n\delta^2}
\to 0$.

$h_n(x) = x^n(1 - x^n)$~: limite simple $0$ sur $\intcc{0}{1}$
(les deux facteurs~; en $x = 1$, $h_n = 0$). Sup~: avec $u = x^n \in
\intcc{0}{1}$, $u(1-u) \leq \frac14$ atteint en $u = \frac12$,
c.-à-d.\ $x = 2^{-1/n} \in \intoo{0}{1}$~: $\norm{h_n}_\infty =
\frac14 \not\to 0$~: pas uniforme sur $\intcc{0}{1}$~; uniforme sur
$\intcc{0}{a}$ ($\sup \leq a^n \to 0$).
\end{solution}

\begin{solution}{exo:b2:funcseq:2}
$\int_0^1 nx\,\eu^{-nx^2}\dd x = \bigl[-\tfrac12
\eu^{-nx^2}\bigr]_0^1 = \frac{1 - \eu^{-n}}{2} \to \frac12$, tandis
que $\int_0^1 \lim g_n = 0$. Pas de contradiction~: le
\cref{thm:b2:funcseq:integration} exige la convergence \emph{uniforme}
sur le segment, qui fait défaut ici (la bosse de hauteur
$\sim\sqrt n$ glisse vers $0$).
\end{solution}

\begin{solution}{exo:b2:funcseq:3}
\omterm{def:b2:funcseq:series}{Convergence normale} sur $\intcc{-1}{1}$~: $\norm{x^n/n^2}_\infty =
\frac{1}{n^2}$, \omterm{def:b2:series:summable}{sommable}~: $S$ y est \omterm{def:b2:metric:continuity}{continue}
(\cref{thm:b2:funcseq:seriestransfer}).

Dérivée~: la série dérivée $\sum \frac{x^{n-1}}{n}$ \omterm{def:b2:integration:improper}{converge}
\omterm{def:b2:funcseq:series}{normalement} sur tout $\intcc{-a}{a}$, $a < 1$ ($\sup =
\frac{a^{n-1}}{n}$)~: $S$ est de classe $C^1$ sur $\intoo{-1}{1}$ avec
\[
S'(x) = \sum_{n\geq1} \frac{x^{n-1}}{n} = \frac1x \sum_{n\geq1}
\frac{x^n}{n} = -\frac{\ln(1 - x)}{x}
\qquad (0 < \abs x < 1),
\]
la dernière identité étant la série logarithmique de première année
(redémontrée honnêtement au \cref{ch:b2:powerseries}).
\end{solution}

\begin{solution}{exo:b2:funcseq:4}
Pour $x > 0$ fixé la série est \omterm{def:b2:linalg:alternating}{alternée} avec $\frac{1}{n + x}
\downarrow 0$~: \omterm{def:b2:funcseq:def}{convergence simple}, et la majoration du reste
$\abs{R_N(x)} \leq \frac{1}{N + 1 + x} \leq \frac{1}{N+1}$ est
\emph{uniforme} sur $\intco{\delta}{\infty}$ (et même sur
$\intoo{0}{\infty}$)~: \omterm{def:b2:funcseq:def}{convergence uniforme}. (Pas normale~:
$\norm{\frac{(-1)^n}{n+x}}_\infty = \frac{1}{n + \delta}$,
divergente.) La continuité découle du
\cref{thm:b2:funcseq:seriestransfer}.

Équation fonctionnelle~: réindexer $F(x + 1)$ avec $m = n + 1$~:
\[
F(x+1) = \sum_{n\geq0} \frac{(-1)^n}{n + 1 + x}
= \sum_{m\geq1}\frac{(-1)^{m-1}}{m+x} ,
\]
donc, en isolant le terme $m = 0$ de $F(x)$,
\[
F(x) + F(x+1)
= \frac{1}{x} + \sum_{m\geq1}
\frac{(-1)^m + (-1)^{m-1}}{m+x} = \frac1x .
\]
\end{solution}

\begin{solution}{exo:b2:funcseq:5}
Posons $g_n = f_n - f \geq 0$ (décroissante en $n$, par hypothèse~; la
limite est $0$ \omterm{def:b2:funcseq:def}{simplement})~; chaque $g_n$ est \omterm{def:b2:metric:continuity}{continue}. Fixons
$\varepsilon > 0$ et posons $U_n = \{x : g_n(x) < \varepsilon\}$~:
\omterm{def:b2:metric:topology}{ouvert} (image réciproque d'un \omterm{def:b2:metric:topology}{ouvert}), croissant ($g_{n+1} \leq g_n$),
et recouvrant $K$ (\omterm{def:b2:funcseq:def}{convergence simple}). Par Borel--Lebesgue
(\cref{thm:b2:metric:borellebesgue}), un nombre fini de $U_{n_1}
\subseteq \dots \subseteq U_{n_k}$ recouvrent $K$~: donc $K = U_{n_k}$,
c.-à-d.\ $\norm{g_{n_k}}_\infty \leq \varepsilon$, et par monotonie
$\norm{g_n}_\infty \leq \varepsilon$ pour tout $n \geq n_k$~:
\omterm{def:b2:funcseq:def}{convergence uniforme}. (La monotonie est essentielle~: les bosses
glissantes de l'\cref{exo:b2:funcseq:2} convergent \omterm{def:b2:funcseq:def}{simplement} sur un
\omterm{def:b2:metric:compact}{compact} sans uniformité.)
\end{solution}

\begin{solution}{exo:b2:funcseq:6}
Substituer $u = nx$~:
\[
\int_0^1 \frac{n f(x)}{1 + n^2x^2}\dd x
= \int_0^n \frac{f(u/n)}{1 + u^2}\,\dd u .
\]
Les intégrandes $h_n(u) = \frac{f(u/n)}{1+u^2}\mathbf{1}_{u \leq n}$
convergent \omterm{def:b2:funcseq:def}{simplement} vers $\frac{f(0)}{1+u^2}$ (continuité de $f$ en
$0$) et sont dominées par $\frac{\norm f_\infty}{1 + u^2}$,
intégrable sur $\intco{0}{\infty}$~: la convergence dominée
(\cref{thm:b2:integration:dominated}) donne la limite
\[
\int_0^\infty \frac{f(0)}{1 + u^2}\dd u = \frac{\pi}{2} f(0) .
\]
(Les noyaux se concentrent en $0$~: une approximation de l'identité.)
\end{solution}

\begin{solution}{exo:b2:funcseq:7}
Pour $f(x) = x^2$, utiliser la deuxième famille d'identités
binomiales de la preuve~: $\sum_k k^2 p_k = n(n-1)x^2 + nx$. D'où
\[
B_n(f)(x) = \sum_k \frac{k^2}{n^2}\,p_k
= \frac{n(n-1)x^2 + nx}{n^2}
= x^2 + \frac{x(1 - x)}{n} :
\]
$B_n(f) - f = \frac{x(1-x)}{n}$, de \omterm{def:b2:nvs:norm}{norme} sup $\frac{1}{4n} =
O\bigl(\frac1n\bigr)$, comme prévu.
\end{solution}

\begin{solution}{exo:b2:funcseq:8}
Pour $\varepsilon = 1$ il existe $N$ tel que $\norm{P_n - P_m}_{\infty,
\R} \leq 1$ pour $m, n \geq N$. Un polynôme borné sur $\R$ est
constant (un polynôme non constant tend vers $\pm\infty$)~: $P_n - P_m =
c_{n,m}$, constantes. Donc pour $n \geq N$~: $P_n = P_N + c_n$ avec
$c_n = P_n(0) - P_N(0)$ convergente (\omterm{def:b2:funcseq:def}{convergence simple} en $0$).
Ainsi $f = \lim P_n = P_N + \lim c_n$~: un polynôme.
\end{solution}

\begin{solution}{exo:b2:funcseq:9}
(a) $\norm{(3/4)^n\varphi(4^n\cdot)}_\infty = \frac12
(3/4)^n$~: \omterm{def:b2:funcseq:series}{convergence normale}, donc $W$ est \omterm{def:b2:metric:continuity}{continue}
(\cref{thm:b2:funcseq:seriestransfer}).

(b) Fixons $x$, $m$~; choisissons le signe de $h_m = \pm\frac12 4^{-m}$
de sorte que le segment $\intcc{4^mx}{4^m(x + h_m)}$ (de longueur
$\frac12$) ne contienne aucun demi-entier, rendant $\varphi$ affine de
pente $\pm1$ sur celui-ci (possible~: un intervalle de longueur
$\frac12$ rencontre au plus un point demi-entier~; choisir le côté qui
l'évite).

Pour $n > m$~: $4^n h_m = \pm\frac12 4^{n-m}$ est un entier, et
$\varphi$ est $1$-périodique~: le $n$-ième terme de la différence
s'annule.

Pour $n = m$~: $\abs{\varphi(4^m x + 4^m h_m) - \varphi(4^m x)} =
\abs{4^m h_m} = \frac12$ ($\varphi$ affine de pente $\pm 1$ sur le
segment), donc le terme contribue exactement $(3/4)^m \cdot
\frac{1/2}{\abs{h_m}} = (3/4)^m\,4^m = 3^m$ dans le quotient.

Pour $n < m$~: le $\varphi$ $1$-lipschitzien donne
$\bigl|(3/4)^n\bigl(\varphi(4^nx + 4^nh_m) -
\varphi(4^nx)\bigr)\bigr| \leq (3/4)^n 4^n\abs{h_m} =
3^n\abs{h_m}$~: chacun contribue au plus $3^n$ dans le quotient.

D'où
\[
\Bigl|\frac{W(x + h_m) - W(x)}{h_m}\Bigr|
\geq 3^m - \sum_{n=0}^{m-1} 3^n
= 3^m - \frac{3^m - 1}{2} = \frac{3^m + 1}{2}
\longrightarrow \infty .
\]
Si $W$ était dérivable en $x$, tout taux d'accroissement le long de
$h_m \to 0$ convergerait vers $W'(x)$~: contradiction. $W$ est
\omterm{def:b2:metric:continuity}{continue} partout, dérivable nulle part.
\end{solution}

\begin{solution}{exo:b2:funcseq:10}
\omterm{def:b2:funcseq:def}{Simplement}~: pour $x \in \intco{0}{1}$ la série est géométrique de
raison $-x$,
\[
\sum_{n\geq0}(-1)^n x^n(1-x) = \frac{1-x}{1+x},
\]
et en $x = 1$ tout terme s'annule~: somme $0 = \frac{1-1}{2}$,
cohérent --- la somme est \omterm{def:b2:metric:continuity}{continue} sur $\intcc{0}{1}$.
Uniformité~: le reste est une queue géométrique,
\[
\abs{R_N(x)} = \frac{x^{N+1}(1-x)}{1+x} \leq x^{N+1}(1 - x)
\leq \max_{\intcc01} t^{N+1}(1-t)
= \frac{1}{N+2}\Bigl(\frac{N+1}{N+2}\Bigr)^{\!N+1}
\leq \frac{1}{N+2} \to 0 ,
\]
\omterm{def:b2:funcseq:def}{uniformément} en $x$. Pas normale~: $\norm{u_n}_\infty =
\max x^n(1-x) = \frac{1}{n+1}\bigl(\frac{n}{n+1}\bigr)^n \sim
\frac{1}{\eu\,n}$, et $\sum \frac1{\eu n}$ diverge. Contraste~:
$\sum x^n(1-x)$ a pour sommes partielles $1 - x^{N+1}$, convergeant
\omterm{def:b2:funcseq:def}{simplement} vers la fonction \emph{discontinue} $\mathbf 1_{\intco01}$~:
d'après le \cref{thm:b2:funcseq:continuity}, cette convergence ne peut
être uniforme sur $\intcc{0}{1}$.
\end{solution}

\begin{solution}{exo:b2:funcseq:11}
La limite $f$ est \omterm{def:b2:metric:continuity}{continue} (\cref{thm:b2:funcseq:continuity}).
Alors
\[
\abs{f_n(x_n) - f(x)}
\leq \abs{f_n(x_n) - f(x_n)} + \abs{f(x_n) - f(x)}
\leq \norm{f_n - f}_\infty + \abs{f(x_n) - f(x)} ,
\]
et les deux termes tendent vers $0$ (\omterm{def:b2:funcseq:def}{convergence uniforme}~; continuité
de $f$ en $x$). Contre-exemple sous simple convergence~:
$g_n(x) = nx\,\eu^{-nx^2} \to 0$ \omterm{def:b2:funcseq:def}{simplement} sur $\intcc{0}{1}$
avec $g_n$ et la limite \omterm{def:b2:metric:continuity}{continues}, pourtant en $x_n =
\frac{1}{\sqrt{2n}} \to 0$~:
\[
g_n(x_n) = \sqrt{\frac n2}\,\eu^{-1/2} \longrightarrow +\infty
\neq 0 = f(0) .
\]
\end{solution}

\begin{solution}{exo:b2:funcseq:12}
\begin{enumerate}
  \item Récurrence. $n = 1$ est la définition. Supposons la formule
        pour $n$ et posons $g(x) = \int_0^x \frac{(x-t)^n}{n!}
        f(t)\dd t$. Pour une intégrande \omterm{def:b2:metric:continuity}{continue} en $(x,t)$ et
        $C^1$ en $x$, l'\omterm{thm:b2:integration:continuity}{intégrale à paramètre} à borne variable se
        dérive en
        \[
        g'(x) = \frac{(x-x)^n}{n!}f(x)
        + \int_0^x \frac{(x-t)^{n-1}}{(n-1)!}f(t)\dd t
        = T^nf(x)
        \]
        (scinder $g(x+h) - g(x)$ en la bande
        $\int_x^{x+h}$, qui est $O(h\cdot\sup)$ avec l'intégrande
        s'annulant en $t = x$ comme $h^n$, et l'intégrale fixe de
        l'accroissement en $x$, traitée par l'inégalité des
        accroissements finis et la continuité). De plus $(T^{n+1}f)' =
        T^nf$ (théorème fondamental de l'analyse) et $g(0) =
        T^{n+1}f(0) = 0$~: deux primitives de $T^nf$ s'annulant en
        $0$ coïncident, donc $T^{n+1}f = g$. La majoration~:
        \[
        \abs{T^nf(x)} \leq \norm f_\infty
        \int_0^x \frac{(x-t)^{n-1}}{(n-1)!}\dd t
        = \norm f_\infty\,\frac{x^n}{n!}
        \leq \frac{\norm f_\infty}{n!} .
        \]
  \item $\sum_n \norm{T^nf}_\infty \leq \eu\,\norm f_\infty$~:
        \omterm{def:b2:funcseq:series}{convergence normale}, donc uniforme~; $S$ est \omterm{def:b2:metric:continuity}{continue}.
        Les sommes partielles vérifient $S_N = f + T S_{N-1}$, et $T$
        est $1$-lipschitzien pour $\norm\cdot_\infty$
        ($\abs{Tg(x)} \leq x\norm g_\infty$)~: en faisant $N \to
        \infty$ des deux côtés on obtient $S = f + TS$.
  \item Posons $V(x) = f(x) + \eu^x\int_0^x \eu^{-t}f(t)\dd t$.
        Alors $V - f$ est $C^1$ avec $(V-f)'(x) =
        \eu^x\int_0^x\eu^{-t}f + f(x) = V(x)$, et $(TV)' = V$
        avec $(V - f)(0) = TV(0) = 0$~: donc $V - f = TV$, c.-à-d.\
        $V$ résout l'équation. Unicité~: si $S_1, S_2$ sont des
        solutions \omterm{def:b2:metric:continuity}{continues}, $D = S_1 - S_2$ vérifie $D =
        TD$, donc $D = T^nD$ pour tout $n$ et $\norm D_\infty
        \leq \frac{\norm D_\infty}{n!} \to 0$~: $D = 0$.
        Par conséquent
        \[
        \sum_{n\geq0} T^nf(x) = f(x) +
        \int_0^x \eu^{x-t}f(t)\,\dd t .
        \]
        (La série $\sum T^n$ est une série géométrique
        d'opérateurs~: un premier avant-goût de la résolvante
        $(\mathrm{Id} - T)^{-1}$, développée dans le volume de
        troisième année.)
\end{enumerate}
\end{solution}

\begin{solution}{pb:b2:funcseq:1}
\textbf{1.} La linéarité est claire d'après la formule. Positivité~:
les poids $p_k(x) \geq 0$, donc $f \geq 0$ force $B_nf \geq
0$~; la monotonie s'ensuit en l'appliquant à $g - f$. Borne~: $\pm f \leq
\norm f_\infty$ donne $\pm B_nf \leq \norm f_\infty B_ne_0 =
\norm f_\infty$. Extrémités~: $p_k(0) = \mathbf 1_{k=0}$ et
$p_k(1) = \mathbf 1_{k=n}$, donc $B_nf(0) = f(0)$, $B_nf(1) =
f(1)$.

\textbf{2.} Dériver $(x+y)^n = \sum_k\binom nk x^ky^{n-k}$
en $x$, multiplier par $x$, et poser $y = 1 - x$~:
\[
nx = \sum_k k\,p_k(x) ;
\]
deux fois, en multipliant par $x^2$~: $n(n-1)x^2 = \sum_k k(k-1)p_k(x)$.
D'où $B_ne_0 = 1$ (théorème du binôme), $B_ne_1(x) =
\frac{nx}{n} = x$, et
\[
B_ne_2(x) = \frac{\sum_k k^2p_k}{n^2}
= \frac{n(n-1)x^2 + nx}{n^2}
= x^2 + \frac{x(1-x)}{n} .
\]

\textbf{3.} Développer~:
\[
\sum_k\Bigl(\frac kn - x\Bigr)^{\!2} p_k
= B_ne_2(x) - 2x\,B_ne_1(x) + x^2
= \frac{x(1-x)}{n} .
\]
Cauchy--Schwarz avec le découpage $\abs{k/n - x}\sqrt{p_k}
\cdot \sqrt{p_k}$~:
\[
\sum_k\Bigl|\frac kn - x\Bigr| p_k
\leq \Bigl(\sum_k\Bigl(\frac kn - x\Bigr)^2
p_k\Bigr)^{\!1/2}
= \sqrt{\frac{x(1-x)}{n}} \leq \frac{1}{2\sqrt n},
\]
en utilisant $x(1-x) \leq \frac14$.

\textbf{4.} Les poids $p_k(x)$ sont positifs de somme $1$
et de barycentre $\sum_k \frac kn p_k(x) = x$ (question 2). L'inégalité
de Jensen finie pour $f$ convexe (récurrence à partir de la définition
à deux points, volume de première année) donne
\[
f(x) = f\Bigl(\sum_k \frac kn\,p_k\Bigr)
\leq \sum_k f\Bigl(\frac kn\Bigr)p_k = B_nf(x) .
\]

\textbf{5.} Sur $\{k : \abs{k/n - x} > \delta\}$ on a
$\bigl(\frac{k/n - x}{\delta}\bigr)^2 > 1$, donc
\[
\sum_{\abs{k/n-x}>\delta} p_k
\leq \frac{1}{\delta^2}\sum_k\Bigl(\frac kn -
x\Bigr)^2p_k
= \frac{x(1-x)}{n\delta^2} \leq \frac{1}{4n\delta^2} .
\]
Lecture~: $p_k(x)$ est la loi d'une fréquence empirique $S_n/n$ de
$n$ lancers de pièce de biais $x$~; sa moyenne est $x$, sa variance
$\frac{x(1-x)}n \to 0$, et l'affichage est l'inégalité de Tchebychev~:
la masse se concentre en $x$, si bien que moyenner $f$
contre elle reproduit $f(x)$ à la limite
(\cref{ch:b2:genfun} rend le vocabulaire officiel).

\textbf{6.} $\omega \leq 2M < \infty$~; la monotonie est claire
(sup sur un ensemble plus grand). Heine~: $f$ \omterm{def:b2:metric:continuity}{continue} sur un \omterm{def:b2:metric:compact}{compact}
est \omterm{def:b2:funcseq:def}{uniformément} \omterm{def:b2:metric:continuity}{continue}, ce qui dit exactement $\omega(\delta) \to 0$
quand $\delta \to 0^+$. Sous-additivité~: si $\abs{s - t} \leq
\delta_1 + \delta_2$, le point $u$ du segment $\intcc st$
à distance $\min(\delta_1, \abs{s-t})$ de $s$ vérifie
$\abs{s-u} \leq \delta_1$, $\abs{u-t} \leq \delta_2$, et
$\abs{f(s)-f(t)} \leq \abs{f(s)-f(u)} + \abs{f(u)-f(t)}$.
En itérant, $\omega(p\delta) \leq p\,\omega(\delta)$ pour $p \in
\N^*$~; pour $\lambda > 0$, avec $p = \lceil\lambda\rceil \leq 1
+ \lambda$~: $\omega(\lambda\delta) \leq \omega(p\delta) \leq
p\,\omega(\delta) \leq (1+\lambda)\omega(\delta)$.

\textbf{7.} Puisque $\sum p_k = 1$~:
\[
\abs{B_nf(x) - f(x)}
= \Bigl|\sum_k\bigl(f(k/n) - f(x)\bigr)p_k\Bigr|
\leq \sum_k\omega\bigl(\abs{k/n - x}\bigr)p_k .
\]
Pour chaque $k$, la question 6 avec $\lambda = \abs{k/n - x}/\delta$
donne $\omega(\abs{k/n-x}) \leq \bigl(1 +
\frac{\abs{k/n-x}}\delta\bigr)\omega(\delta)$~; en sommant contre
les $p_k$ on obtient l'estimation maîtresse.

\textbf{8.} Insérer la borne de la question 3~:
\[
\abs{B_nf(x) - f(x)} \leq \Bigl(1 +
\frac{1}{2\delta\sqrt n}\Bigr)\omega(\delta),
\]
\omterm{def:b2:funcseq:def}{uniformément} en $x$~; avec $\delta = n^{-1/2}$ la parenthèse vaut
$\frac32$~: $\norm{B_nf - f}_\infty \leq
\frac32\omega(n^{-1/2}) \to 0$ par la question 6 (Heine). Ceci
redémontre le \cref{thm:b2:funcseq:weierstrass} avec une vitesse.

\textbf{9.} $L$-lipschitzienne signifie $\omega(\delta) \leq L\delta$~:
vitesse $\frac{3L}{2\sqrt n}$. $\alpha$-höldérienne signifie
$\omega(\delta) \leq C\delta^\alpha$~: vitesse $\frac{3C}{2}
n^{-\alpha/2}$.

\textbf{10.} En utilisant $k\binom{2m}k = 2m\binom{2m-1}{k-1}$~:
\[
\sum_{k=m+1}^{2m}k\binom{2m}k
= 2m\sum_{j=m}^{2m-1}\binom{2m-1}{j}
= 2m\cdot 2^{2m-2},
\]
car $j \mapsto 2m-1-j$ bijecte $\{m,\dots,2m-1\}$ sur
$\{0,\dots,m-1\}$, donc la somme est la moitié de $2^{2m-1}$. De plus
$\sum_{k=m+1}^{2m}\binom{2m}k = \frac{2^{2m} -
\binom{2m}m}{2}$ (même symétrie). D'où
\[
\sum_{k=m+1}^{2m}(k-m)\binom{2m}k
= m\,2^{2m-1} - m\,\frac{2^{2m} - \binom{2m}m}{2}
= \frac m2\binom{2m}m .
\]
La substitution $k \mapsto 2m-k$ envoie les termes avec $k < m$
sur ceux avec $k > m$ (binomiaux égaux, $\abs{k-m}$ égaux)~:
la somme absolue est le double de la somme unilatérale, $m\binom{2m}m$.

\textbf{11.} En $x = \frac12$, $p_k(\tfrac12) =
\binom{2m}k2^{-2m}$ et $f(\tfrac12) = 0$~:
\[
B_{2m}f\Bigl(\frac12\Bigr)
= \sum_k\Bigl|\frac{k}{2m} - \frac12\Bigr|
\binom{2m}k 2^{-2m}
= \frac{2^{-2m}}{2m}\,m\binom{2m}m
= \frac{\binom{2m}m}{2\cdot4^m}
\sim \frac{1}{2\sqrt{\pi m}}
\]
d'après l'\cref{ex:b2:comparison:centralbinomial}. Puisque $\omega_f
(\delta) = \delta$ ici (la fonction est $1$-lipschitzienne et la
borne est atteinte), la question 8 prédit au plus
$\frac32(2m)^{-1/2}$~: la véritable erreur $\frac{1}{2\sqrt{\pi m}}$
a exactement l'\omterm{def:b2:structures:generated}{ordre} $n^{-1/2}$ --- la vitesse est optimale à la
constante près.

\textbf{12.} $f \leq g$ donne $g - f \geq 0$, donc $L(g-f) \geq
0$, c.-à-d.\ $Lf \leq Lg$. De $-\abs f \leq f \leq \abs f$~:
$-L\abs f \leq Lf \leq L\abs f$, c.-à-d.\ $\abs{Lf} \leq
L\abs f$.

\textbf{13.} Par Heine choisir $\delta$ avec $\abs{f(s)-f(x)}
\leq \varepsilon$ dès que $\abs{s-x} \leq \delta$. Si
$\abs{s - x} > \delta$, alors $\frac{(s-x)^2}{\delta^2} > 1$ et
$\abs{f(s)-f(x)} \leq 2M \leq \frac{2M}{\delta^2}(s-x)^2$. Dans
les deux cas la borne annoncée vaut.

\textbf{14.} Fixer $x$~; la question 13 dit, en tant que fonctions de
$s$~:
\[
-\varepsilon e_0 - \frac{2M}{\delta^2}q_x
\;\leq\; f - f(x)e_0
\;\leq\; \varepsilon e_0 + \frac{2M}{\delta^2}q_x,
\qquad q_x = e_2 - 2x\,e_1 + x^2e_0 .
\]
Appliquer le $L_n$ monotone linéaire (question 12) et évaluer en
$x$~:
\[
\abs{L_nf(x) - f(x)L_ne_0(x)}
\leq \varepsilon L_ne_0(x)
+ \frac{2M}{\delta^2}\bigl(L_ne_2(x) - 2xL_ne_1(x)
+ x^2L_ne_0(x)\bigr) .
\]

\textbf{15.} Écrire $\alpha_j = L_ne_j - e_j$, de sorte que
$\norm{\alpha_j}_\infty \to 0$. Puisque $e_2(x) - 2xe_1(x) +
x^2e_0(x) = 0$~:
\[
L_ne_2(x) - 2xL_ne_1(x) + x^2L_ne_0(x)
= \alpha_2(x) - 2x\,\alpha_1(x) + x^2\alpha_0(x),
\]
de \omterm{def:b2:nvs:norm}{norme} sup au plus $\norm{\alpha_2} + 2\norm{\alpha_1} +
\norm{\alpha_0} \to 0$. De plus $L_ne_0 \to e_0$ \omterm{def:b2:funcseq:def}{uniformément}, donc
$L_ne_0 \leq 2$ pour $n$ grand, et $\abs{f(x)}\abs{L_ne_0(x) -
1} \leq M\norm{\alpha_0} \to 0$. En assemblant avec la question 14~:
pour $n$ grand, \omterm{def:b2:funcseq:def}{uniformément} en $x$,
\[
\abs{L_nf(x) - f(x)} \leq 2\varepsilon +
\frac{2M}{\delta^2}\,o(1) + M\,o(1) \leq 3\varepsilon :
\]
$L_nf \to f$ \omterm{def:b2:funcseq:def}{uniformément} --- le théorème de Korovkin.

\textbf{16.} $B_ne_0 = e_0$ et $B_ne_1 = e_1$ exactement, et
$\norm{B_ne_2 - e_2}_\infty = \max_x\frac{x(1-x)}{n} =
\frac{1}{4n} \to 0$ (question 2)~: Korovkin s'applique, et
Weierstrass s'ensuit pour la troisième fois.

\textbf{17.} $I_nf$ est linéaire en $f$ (les valeurs nodales le sont),
et sur chaque maille l'interpolant affine de valeurs nodales positives
est positif~: positif. $I_ne_0 = e_0$ et $I_ne_1 = e_1$
car une fonction affine égale son propre interpolant. Sur une
maille $\intcc ab$ ($b - a = \frac1n$), l'interpolant affine de
$e_2$ est $L(t) = (a+b)t - ab$, et
\[
L(t) - t^2 = (t-a)(b-t) \in
\intcc{0}{\tfrac{(b-a)^2}{4}} ,
\]
avec le maximum au milieu~: $\norm{I_ne_2 - e_2}_\infty =
\frac{1}{4n^2} \to 0$. Korovkin~: $I_nf \to f$ \omterm{def:b2:funcseq:def}{uniformément} pour
toute $f$ \omterm{def:b2:metric:continuity}{continue} --- approximation polygonale, sans
estimation supplémentaire.

\textbf{18.} Étant donnés $f$ et $\varepsilon$~: Weierstrass fournit
un polynôme $P = \sum_{j=0}^d a_jx^j$ avec $\norm{f -
P}_\infty \leq \frac\varepsilon2$~; remplacer chaque $a_j$ par un
rationnel $b_j$ avec $\abs{a_j - b_j} \leq
\frac{\varepsilon}{2(d+1)}$ déplace la \omterm{def:b2:nvs:norm}{norme} sup sur
$\intcc{0}{1}$ d'au plus $\frac\varepsilon2$. L'ensemble des
polynômes à coefficients rationnels est une réunion \omterm{def:b2:structures:countable}{dénombrable} (sur
$d$) d'\omterm{def:b2:structures:countable}{ensembles dénombrables}, donc \omterm{def:b2:structures:countable}{dénombrable}, et dense~:
$C(\intcc{0}{1})$ est séparable.

\textbf{19.} Par linéarité $\int_0^1 fP = 0$ pour tout
polynôme $P$. Choisir des polynômes $P_n \to f$ \omterm{def:b2:funcseq:def}{uniformément}
(Weierstrass)~:
\[
\Bigl|\int_0^1 f^2\Bigr|
= \Bigl|\int_0^1 f\,(f - P_n)\Bigr|
\leq \norm f_\infty\,\norm{f - P_n}_\infty
\longrightarrow 0 ,
\]
donc $\int_0^1 f^2 = 0$. Si $f(x_0) \neq 0$, la continuité donne
$f^2 \geq c > 0$ sur un sous-intervalle, contredisant l'intégrale
nulle~: $f = 0$. Par conséquent deux fonctions \omterm{def:b2:metric:continuity}{continues} de
mêmes moments $\int f t^n$ coïncident.

\textbf{20.} Avec $p_{n,k}(x) = \binom nk x^k(1-x)^{n-k}$ et
les conventions $p_{n-1,-1} = p_{n-1,n} = 0$, la règle du produit
et $k\binom nk = n\binom{n-1}{k-1}$, $(n-k)\binom nk =
n\binom{n-1}{k}$ donnent
\[
p_{n,k}'(x) = n\bigl(p_{n-1,k-1}(x) - p_{n-1,k}(x)\bigr) .
\]
En sommant contre $f(k/n)$ et en décalant l'indice dans la première
somme (\omterm{thm:b2:series:abel}{transformation d'Abel})~:
\[
(B_nf)'(x) = n\sum_{j=0}^{n-1}\Bigl(f\Bigl(\frac{j+1}n\Bigr)
- f\Bigl(\frac jn\Bigr)\Bigr)p_{n-1,j}(x) .
\]

\textbf{21.} Par le théorème des accroissements finis, $f(\frac{j+1}n) -
f(\frac jn) = \frac1n f'(\xi_j)$ avec $\xi_j \in
\intoo{j/n}{(j+1)/n}$, donc $(B_nf)'(x) = \sum_j
f'(\xi_j)\,p_{n-1,j}(x)$. Le nœud $\frac{j}{n-1}$ appartient aussi à
$\intcc{j/n}{(j+1)/n}$ (les deux inégalités se réduisent à $j \leq
n-1$), donc $\abs{\xi_j - \frac j{n-1}} \leq \frac1n$ et
\[
\bigl|(B_nf)'(x) - B_{n-1}(f')(x)\bigr|
\leq \sum_j\Bigl|f'(\xi_j) -
f'\Bigl(\frac{j}{n-1}\Bigr)\Bigr| p_{n-1,j}(x)
\leq \omega_{f'}\Bigl(\frac1n\Bigr) \longrightarrow 0
\]
\omterm{def:b2:funcseq:def}{uniformément}. Puisque $B_{n-1}(f') \to f'$ \omterm{def:b2:funcseq:def}{uniformément}
(\cref{thm:b2:funcseq:weierstrass} appliqué au $f'$ \omterm{def:b2:metric:continuity}{continu}),
l'inégalité triangulaire donne $(B_nf)' \to f'$
\omterm{def:b2:funcseq:def}{uniformément}. Les polynômes $P_n = B_nf$ convergent alors vers $f$
au sens $C^1$.

\textbf{22.} Taylor--Lagrange en $x$~: $f(\frac kn) - f(x) =
f'(x)(\frac kn - x) + \frac{f''(\xi_k)}2(\frac kn - x)^2$.
En sommant contre $p_k$, le terme linéaire meurt (question 2)~:
\[
\abs{B_nf(x) - f(x)}
\leq \frac{\norm{f''}_\infty}{2}\sum_k\Bigl(\frac kn -
x\Bigr)^2p_k
= \frac{\norm{f''}_\infty}{2}\cdot\frac{x(1-x)}{n}
\leq \frac{\norm{f''}_\infty}{8n} .
\]

\textbf{23.} Deux dérivations de plus de $(x+y)^n$ donnent les
moments factoriels, avec $n_{(j)} = n(n-1)\cdots(n-j+1)$~:
\[
\sum_k k_{(j)}\,p_k = n_{(j)}\,x^j \qquad (j = 3, 4),
\]
et $k^3 = k_{(3)} + 3k_{(2)} + k$, $k^4 = k_{(4)} + 6k_{(3)} +
7k_{(2)} + k$ les convertissent en moments de puissances~:
\[
\sum_k k^3p_k = n_{(3)}x^3 + 3n_{(2)}x^2 + nx,
\qquad
\sum_k k^4p_k = n_{(4)}x^4 + 6n_{(3)}x^3 + 7n_{(2)}x^2 + nx .
\]
En développant $(k - nx)^4$ et en collectant (un calcul patient mais
purement mécanique avec les quatre moments de puissances)~:
\[
\sum_k(k-nx)^4p_k = nx(1-x)\bigl(1 + 3(n-2)x(1-x)\bigr) .
\]
Avec $x(1-x) \leq \frac14$~: le membre de droite est au plus
$\frac n4\bigl(1 + \frac{3n}4\bigr) = \frac{3n^2}{16} +
\frac n4 \leq n^2$ pour $n \geq 1$.

\textbf{24.} La forme de Peano de Taylor en $x$ définit $\eta(t) =
\frac{f(t) - f(x) - f'(x)(t-x) - \frac12f''(x)(t-x)^2}
{(t-x)^2}$ pour $t \neq x$, $\eta(x) = 0$~: par Taylor--Lagrange
$\eta(t) = \frac12\bigl(f''(\xi) - f''(x)\bigr)$ pour un certain $\xi$
entre $t$ et $x$, donc $\abs\eta \leq \norm{f''}_\infty$ et
$\eta(t) \to 0$ quand $t \to x$ (continuité de $f''$). En sommant
le développement contre $p_k$ et en utilisant les questions 2--3~:
\[
n\bigl(B_nf(x) - f(x)\bigr)
= \frac{x(1-x)}{2}f''(x)
+ n\sum_k\eta\Bigl(\frac kn\Bigr)\Bigl(\frac kn -
x\Bigr)^2p_k .
\]
Étant donné $\varepsilon$, choisir $\delta$ avec $\abs\eta \leq
\varepsilon$ sur $\abs{t - x}\leq\delta$. Partie proche~: au plus
$\varepsilon\,n\cdot\frac{x(1-x)}n \leq \varepsilon$. Partie lointaine~:
avec $C = \norm{f''}_\infty$ et la question 23,
\[
n\,C\sum_{\abs{k/n-x}>\delta}\Bigl(\frac kn -
x\Bigr)^2p_k
\leq \frac{nC}{\delta^2}\sum_k\Bigl(\frac kn -
x\Bigr)^4p_k
= \frac{nC}{\delta^2 n^4}\sum_k(k-nx)^4p_k
\leq \frac{C}{\delta^2 n} \longrightarrow 0 .
\]
D'où $n(B_nf(x) - f(x)) \to \frac{x(1-x)}2f''(x)$ ---
le théorème de Voronovskaya. Pour $f = e_2$ c'est exact à tout
$n$ (\cref{exo:b2:funcseq:7})~: le plafond en $\frac1n$ est réel.

\textbf{25.} (i) La positivité a transformé des inégalités
ponctuelles en inégalités d'opérateurs~: elle a donné la borne de
\omterm{def:b2:nvs:norm}{norme}, Jensen, Tchebychev, et tout Korovkin --- la linéarité seule ne
prouve rien ici. (ii) La compacité est intervenue par Heine (questions
6, 13), par la bornitude de $f$, et par la \omterm{def:b2:nvs:norm}{norme} même
$\norm\cdot_\infty$ qui est finie. (iii) Trois fonctions test
suffisent car la positivité réduit tout au contrôle de
$L_n$ sur l'unique famille $(s-x)^2 = e_2 - 2xe_1 + x^2e_0$,
dont l'espace engendré est celui de $e_0, e_1, e_2$. (iv) Bernstein
\omterm{def:b2:integration:improper}{converge} à la vitesse honnête $\omega(n^{-1/2})$ pour toute $f$
\omterm{def:b2:metric:continuity}{continue} (optimale, question 11) mais sature à $\frac1n$ pour $f$
régulière (question 24)~; le chapitre de Fourier mène le même
programme pour les fonctions périodiques avec le noyau de Fejér ---
un autre opérateur positif aux mêmes vertus et à la même modestie.
\end{solution}
